%

\section{Abordagem do Método da Homogeneização Assintótica (MHA)}

\begin{figure}[!h]
    \centering
    \includegraphics[width=0.5\linewidth]{images/Modelo_do_compósito.png}
    \caption{Representation of the full micro-periodic composite pressure vessel, consisting of two phases of individually homogeneous constituents and radius $r_i$ ($i = 1, 2, \dots, N$).}
    \label{fig:Composito}
\end{figure}

In this section, the general form of the \textcolor{blue}{homogeneous Cauchy momentum equation} \cite{sadd2009elasticity} is taken as the starting point to analyze the mechanical behavior of a cylindrical pressure vessel composed of periodically alternating material layers in the radial direction, as illustrated in Fig.~\ref{fig:Composito}. The tube is subjected to a constant internal pressure. All constituent materials are assumed to exhibit a well-defined linear elastic regime, remaining within this regime throughout the process. The governing equations are given by:
%
\begin{eqnarray}\label{eq:OriginalPDE}
    -\nabla\cdot\boldsymbol{\sigma}^{\varepsilon} = 0, \quad \mathbf{x} \in \Omega, \nonumber\\
    \boldsymbol{\sigma}^{\varepsilon} = \lambda^{\varepsilon} \, \mathrm{tr}(\boldsymbol{\epsilon}^{\varepsilon}) \, I + 2\mu^{\varepsilon} \boldsymbol{\epsilon}^{\varepsilon}, \quad
    \boldsymbol{\epsilon}^{\varepsilon} = \frac{1}{2}\left[\nabla \mathbf{u}^{\varepsilon} + (\nabla \mathbf{u}^{\varepsilon})^T\right], \\
    \boldsymbol{\sigma}^{\varepsilon} \cdot \mathbf{n} = -\mathbf{p}, \quad \mathbf{x} \in \Gamma_u, \quad
    \boldsymbol{\sigma}^{\varepsilon} \cdot \mathbf{n} = \mathbf{0}, \quad \mathbf{x} \in \Gamma_t, \nonumber
\end{eqnarray}
%
\noindent where $\boldsymbol{\sigma}^{\varepsilon}$ is the stress tensor, $\mathbf{f}^{\varepsilon}$ denotes the body force field, $\mathbf{u}^{\varepsilon}$ is the displacement vector—the primary unknown—$\mathbf{n}$ is the outward unit normal to the boundary $\partial\Omega$, and $\mathbf{p}$ is the \textcolor{blue}{gauge pressure} applied on the inner surface of the cylinder. The stress–strain relation follows Hooke’s law, $\boldsymbol{\sigma}^{\varepsilon} = \mathcal{\boldsymbol{C}}^{\varepsilon} : \boldsymbol{\epsilon}^{\varepsilon}$, where $\mathcal{\boldsymbol{C}}^{\varepsilon}$ is the elasticity tensor containing the Lamé coefficients $\lambda^{\varepsilon}$ and $\mu^{\varepsilon}$, and $\boldsymbol{\epsilon}^{\varepsilon}$ is the strain tensor. The set $\Gamma_u$ represents the inner surface where constant-pressure Neumann boundary conditions are imposed, while $\Gamma_t$ denotes the traction-free outer surface (see Fig.~\ref{fig:Composito}).

For the case of a cylindrical coordinate system, the first equation in \eqref{eq:OriginalPDE} can be expressed as:
%
\begin{eqnarray}
    \label{eq:EquilibriumCylindricalForm}
    \frac{d\sigma_{r}^{\varepsilon}(r)}{dr} + \frac{1}{r}\left[\sigma_r^{\varepsilon}(r) - \sigma_\theta^{\varepsilon}(r)\right] = 0,
\end{eqnarray}
%
\noindent where $\sigma_r^{\varepsilon}$ and $\sigma_{\theta}^{\varepsilon}$ are the radial and circumferential (hoop) stress components, respectively. The constitutive relations in cylindrical coordinates are:
%
\begin{eqnarray}\label{eq:HookesLawCylindricalCoords}
    \sigma_{r}^{\varepsilon}(r) &=& \left[ \lambda^{\varepsilon}(r) + 2\mu^{\varepsilon}(r) \right] \frac{d u_{r}^{\varepsilon}}{dr} + \lambda^{\varepsilon}(r) \frac{u_r^{\varepsilon}}{r}, \nonumber\\
    \sigma_{\theta}^{\varepsilon}(r) &=& \lambda^{\varepsilon}(r) \frac{d u_{r}^{\varepsilon}}{dr} + \left[ \lambda^{\varepsilon}(r) + 2\mu^{\varepsilon}(r) \right] \frac{u_r^{\varepsilon}}{r}.
\end{eqnarray}

Or, in matrix notation:
%
\begin{equation}
    \begin{aligned}
        \begin{bmatrix}
            \sigma_{r}^{\varepsilon}(r)\\
            \sigma_{\theta}^{\varepsilon}(r)
        \end{bmatrix} =
        \begin{bmatrix}
            \lambda^{\varepsilon}(r) + 2\mu^{\varepsilon}(r) &  \lambda^{\varepsilon}(r)\\
            \lambda^{\varepsilon}(r) & \lambda^{\varepsilon}(r) + 2\mu^{\varepsilon}(r)
        \end{bmatrix}
        \begin{bmatrix}
            \dfrac{du_r^{\varepsilon}}{dr}\\
            \dfrac{u_r^{\varepsilon}}{r}
        \end{bmatrix}.
    \end{aligned}
\end{equation}

\def\Sigmari{\sigma_{r}^{\varepsilon}(r_i) =  C^{\varepsilon}_{11}(r_i)\frac{d u_{r}}{d r} + C^{\varepsilon}_{12}(r_i)\frac{u_r}{r_i}}
\def\Sigmare{\sigma_{r}^{\varepsilon}(r_e) =  C^{\varepsilon}_{11}(r_e)\frac{d u_{r}}{d r} + C^{\varepsilon}_{12}(r_e)\frac{u_r}{r_e}}

\noindent Here, $\mu^{\varepsilon}(r)$ is the shear modulus and $\lambda^{\varepsilon}(r)$ is Lamé’s first parameter. The superscript $\varepsilon$ denotes the scale separation in the material description. By combining the equilibrium equation \eqref{eq:EquilibriumCylindricalForm} with Hooke’s law \eqref{eq:HookesLawCylindricalCoords}, we obtain the differential equation to which the homogenization method will be applied:
%
\def\Ceff{C_{11}^{\varepsilon}(r)}
\def\duh{\frac{du^{\varepsilon}(r)}{dr}}
\def\duhr{\frac{d^2u^{\varepsilon}(r)}{dr^2}}
%
\begin{equation}
\label{eq:EquililbriumFinalVersion}
    \begin{aligned}
        \Ceff \, \duhr + \left( \frac{1}{r} \Ceff + \frac{d\Ceff}{dr} \right) \duh +\\ \left( \frac{1}{r} \frac{d C_{12}^\varepsilon(r)}{dr} - \frac{C_{11}^\varepsilon(r)}{r^2} \right) u^\varepsilon(r) - f^{\varepsilon}(r) = 0,\\
    \Sigmari = -p,\\
    \Sigmare = 0.
    \end{aligned}
\end{equation}

\noindent The coefficients $C^{\varepsilon}_{11}(r) = \lambda^{\varepsilon}(r) + 2\mu^{\varepsilon}(r)$ and $C_{12}^{\varepsilon}(r) = \lambda^{\varepsilon}(r)$, expressed in Voigt notation \cite{Sadd2009}, are the isotropic elasticity tensor components used to simplify Eq.~\eqref{eq:HookesLawCylindricalCoords}. Equation~\eqref{eq:EquililbriumFinalVersion} is subject to the same boundary conditions as Eq.~\eqref{eq:OriginalPDE}.

\section*{Homogenization}

%Definição de variáveis

\newcommand{\stressSAF}[1]{\sigma^{(0)}_#1(r,y) + \varepsilon \sigma^{(1)}_#1(r,y)}
\newcommand{\totalderivative}[2]{\frac{d #1}{d #2}}
\newcommand{\partialderivative}[2]{\frac{\partial #1}{\partial #2}}
\newcommand{\StressSAF}[1]{
\sigma_{#1}^{(0)} + 
\varepsilon\sigma_{#1}^{(1)} + \varepsilon^2\sigma_{#1}^{(2)}
}
\newcommand{\DisplacementSaF}[1]{u_r^{(0)} + \varepsilon u_r^{(1)} + \varepsilon^2 u_r^{(2)}}
\newcommand{\PartialDerivative}[2]{\frac{\partial #1}{\partial #2}}
\newcommand{\SAFAfterChainRule}[3]{\partialderivative{#1^{(0)}}{#2} + 
\varepsilon \partialderivative{#1^{(1)}}{#2} + \varepsilon^2 \partialderivative{#1^{(2)}}{#2} + 
\frac{1}{\varepsilon} \partialderivative{#1^{(0)}}{#3} + \partialderivative{#1^{(1)}}{#3} + 
\varepsilon \partialderivative{#1^{(2)}}{#3}}
\newcommand{\ChainRule}{\partialderivative{}{r} + \frac{1}{\varepsilon} \partialderivative{}{y}}
\newcommand{\DispDerivative}[2]{ \partialderivative{u_r^{(#1)}}{#2} }
\newcommand{\LocalFunctions}{N_{11}(y) \DispDerivative{0}{r} + N_{12}(y)\frac{u_r^{(0)}}{r}}
\newcommand{\DeformationTheta}{\frac{u_r^{(0)}}{r}}
%Stresses
\newcommand{\StressROrderZero}{ \hat{C}_{11} \partialderivative{u_r^{(0)}}{r} + \hat{C}_{12}\DeformationTheta }
\newcommand{\StressThetaOrderZero}{ \hat{C}_{12} \partialderivative{u_r^{(0)}}{r} + \hat{C}_{11}\DeformationTheta}

The asymptotic homogenization method employs a \textcolor{blue}{formal asymptotic solution (FAS)} for the rapidly oscillating boundary value problem (BVP), expressed as:

\begin{equation}\label{eq:asymptotic_expansions}
    \begin{aligned}
        \mathbf{u}^{(\infty)}(r,\varepsilon)=\mathbf{u}^{(0)}(r)+\varepsilon \mathbf{u}^{(1)} (r,y)+\varepsilon^2 \mathbf{u}^{(2)}(r,y) + ...,\\
        \boldsymbol{\epsilon}^{(\infty)}(r,\varepsilon) = \boldsymbol{\epsilon}^{(0)}(r,y) + \varepsilon \boldsymbol{\epsilon}^{(1)}(r,y) +\varepsilon^2 \boldsymbol{\epsilon}^{(2)}(r,y) ...,\\
        \boldsymbol{\sigma}^{(\infty)}(r, \varepsilon) = \boldsymbol{\sigma}^{(0)}(r,y) + \varepsilon \boldsymbol{\sigma}^{(1)}(r,y) +\varepsilon^2 \boldsymbol{\sigma}^{(2)}(r,y) + ...
    \end{aligned}
\end{equation}

Truncating the stress expansion after the first-order term ($\mathcal{O}(\varepsilon)$) in \eqref{eq:asymptotic_expansions} and substituting into the constitutive law \eqref{eq:HookesLawCylindricalCoords} yields the following matrix equation:

\begin{equation}\label{eq:HookesLawWithSAF_MatrixForm}
    \begin{aligned}
        \begin{bmatrix}
            \sigma_{r}^{\varepsilon}(r)\\
            \sigma_{\theta}^{\varepsilon}(r)
        \end{bmatrix} = 
        \begin{bmatrix}
            C_{11}(y) &  C_{12}(y) \\
            C_{12}(y)  & C_{11}(y) 
        \end{bmatrix}
        \begin{bmatrix}
            \frac{du_r^{\varepsilon}}{dr}\\
            \frac{u_r^{\varepsilon}}{r}
        \end{bmatrix}, \\
        \begin{bmatrix}
            \sigma_{r}^{(0)}(r,y) + \varepsilon\sigma_{r}^{(1)}(r,y) + \varepsilon^2\sigma_{r}^{(2)}(r,y)\\
            \sigma_{\theta}^{(0)}(r,y) + \varepsilon\sigma_{\theta}^{(1)}(r,y)+\varepsilon^2\sigma_{\theta}^{(2)}(r,y)
        \end{bmatrix}= \\ 
        \begin{bmatrix}
            C_{11}(y) &  C_{12}(y) \\
            C_{12}(y)  & C_{11}(y) 
        \end{bmatrix}
        \begin{bmatrix}
            \frac{d}{dr}(u_r^{(0)} + \varepsilon u_r^{(1)} + \varepsilon^2 u_r^{(2)})\\
            \frac{1}{r}(u_r^{(0)} + \varepsilon u_r^{(1)} + \varepsilon^2 u_r^{(2)})
        \end{bmatrix}. \\
    \end{aligned}
\end{equation}

Also, by applying the chain rule to Eq. (\ref{eq:HookesLawWithSAF_MatrixForm}), it becomes:

\begin{equation}
    \begin{aligned}
        \begin{bmatrix}
            \StressSAF{r}\\
            \StressSAF{\theta}
        \end{bmatrix} = 
        \begin{bmatrix}
            C_{11}(y) &  C_{12}(y) \\
            C_{12}(y)  & C_{11}(y) 
        \end{bmatrix}
        \begin{bmatrix}
            \left(\frac{\partial}{\partial r} + \frac{1}{\varepsilon}\frac{\partial}{\partial y}\right)(\DisplacementSaF{})\\
            \frac{1}{r}(\DisplacementSaF{})
        \end{bmatrix}. \\
    \end{aligned}
\end{equation}

Expanding the chain rule for the stress-strain relationship yields:

\begin{equation}
    \begin{aligned}
        \begin{bmatrix}
            \StressSAF{r}\\
            \StressSAF{\theta}
        \end{bmatrix} = \\
        \begin{bmatrix}
            C_{11}(y) &  C_{12}(y) \\
            C_{12}(y)  & C_{11}(y) 
        \end{bmatrix}
        \begin{bmatrix}
             \frac{1}{\varepsilon} \frac{\partial u_r^{(0)}}{\partial y} + \frac{\partial u_r^{(0)}}{\partial r} + \frac{\partial u_r^{(1)}}{\partial y} + \varepsilon \frac{\partial u_r^{(1)}}{\partial r}\\
            \frac{1}{r}u_r^{(0)} + \frac{1}{r}\varepsilon u_r^{(1)} + \frac{1}{r}\varepsilon^2 u_r^{(2)} 
        \end{bmatrix}. \\
    \end{aligned}
\end{equation}

Expanding again each stress term, r direction gives:


\begin{equation}
    \begin{aligned} \label{eq:SAFAndHookeInRDirection}
       \StressSAF{r} =\\ C_{11}(y)\left(
             \SAFAfterChainRule{u_r}{r}{y}
        \right) +\\ C_{12}(y)\left(
            \frac{1}{r}u_r^{(0)} + \frac{1}{r}\varepsilon u_r^{(1)} + \frac{1}{r}\varepsilon^2 u_r^{(2)}
        \right).
    \end{aligned}
\end{equation}

Terms with the same power of $\varepsilon$ must be equal, which implies that:

\begin{equation}
    \begin{aligned} \label{eq:HookesLawWithSAFRdirection}
        \varepsilon^{-1}:\partialderivative{u_r^{(0)}}{y} = 0,\\
        \varepsilon^0:\sigma^{(0)}_r = C_{11}(y) \partialderivative{u_r^{(0)}}{r} + C_{11}(y) \partialderivative{u_r^{(1)}}{y} + C_{12}(y)\DeformationTheta, \\
        \varepsilon^1: \sigma^{(1)}_r = C_{11}(y)\partialderivative{u_r^{(1)}}{r} + C_{11}(y)\partialderivative{u_r^{(2)}}{y}+ C_{12}(y)\frac{u_r^{(1)}}{r} , \\
        \varepsilon^2: \sigma^{(2)}_r = C_{11}(y)\partialderivative{u_r^{(2)}}{r} +C_{12}(y) \frac{u_r^{(2)}}{r}.
    \end{aligned}
\end{equation}

Applying these stresses in the equilibrium equation yields:

\begin{equation}
    \begin{aligned}
        \frac{d\sigma_{r}^{\varepsilon}(r)}{dr}+\frac1r( \sigma_r(r)-\sigma_\theta(r))=0,\\
        \left( \ChainRule \right) (\StressSAF{r}) +\\\frac1r [\StressSAF{r} - (\StressSAF{\theta})] = 0.
    \end{aligned}
\end{equation}

Expanding the chain rule:

\begin{equation}
    \begin{aligned}
        \SAFAfterChainRule{\sigma_r}{r}{y} +\\\frac1r [\StressSAF{r} - (\StressSAF{\theta})] = 0.
    \end{aligned}
\end{equation}

The powers of $\varepsilon$ up to the 0th power can again be gathered, as higher powers vanish as $\varepsilon \rightarrow 0$:

\begin{equation}
    \begin{aligned} \label{eq:EquilibriumWithStressSAFs}
        \varepsilon^{-1}: \partialderivative{\sigma_r^{(0)}(r)}{y} = 0,\\
        \varepsilon^{0}: \partialderivative{\sigma_r^{(0)}(r)}{r} + \partialderivative{\sigma_r^{(1)}(r,y)}{y} + \frac{1}{r} \left(  \sigma_r^{(0)}(r) - \sigma_{\theta}^{(0)}(r,y) \right) = 0,
        \end{aligned}
\end{equation}

Expanding the order -1 equation in (\ref{eq:}) using Eq. (\ref{eq:HookesLawWithSAFRdirection}):

\begin{equation}
    \begin{aligned} \label{eq:FindingU1Expresion}
        \partialderivative{}{y} \left( C_{11}(y) \partialderivative{u_r^{(0)}}{r} + C_{11}(y) \partialderivative{u_r^{(1)}}{y} + C_{12}(y)\frac{u_r^{(0)}}{r}\right) = 0 ,\\
        \partialderivative{}{y} \left( C_{11}(y) \partialderivative{u_r^{(0)}}{r} \right) + 
        \partialderivative{}{y} \left( C_{11}(y) \partialderivative{u_r^{(1)}}{r} \right) + 
        \partialderivative{}{y} \left( C_{12}(y)\frac{u_r^{(0)}}{r} \right) = 0,\\
        \partialderivative{}{y} \left( C_{11}(y) \partialderivative{u_r^{(1)}}{r} \right) = -\partialderivative{C_{11}}{y}\DispDerivative{0}{r} - \partialderivative{C_{12}}{y}\frac{u_r^{(0)}}{r}. 
    \end{aligned}
\end{equation}

For convenience, we assume that the first-order displacement function can be written as $ u_r^{(1)} = N_{11}(y) \DispDerivative{0}{r} + N_{12}(y)\frac{u_r^{(0)}}{r} $, with $N_{11}(y)$ and $N_{12}(y)$ being functions defined in the unit cell (local functions). Replacing this expression in Eq. (\ref{eq:FindingU1Expresion}), it yields:

\begin{equation}
    \begin{aligned} \label{eq:PreLocalProblems}
        \partialderivative{}{y} \left( C_{11}(y) \partialderivative{u_r^{(1)}}{r} \right) = -\partialderivative{C_{11}}{y}\DispDerivative{0}{r} - \partialderivative{C_{12}}{y}\frac{u_r^{(0)}}{r},\\
        \partialderivative{}{y}\left[ C_{11}(y) \partialderivative{}{y}\left( \LocalFunctions \right) \right] = -\partialderivative{C_{11}}{y}\DispDerivative{0}{r} - \partialderivative{C_{12}}{y}\frac{u_r^{(0)}}{r},\\
        \partialderivative{}{y} \left( C_{11}(y)\partialderivative{N_{11}(y)}{y}\DispDerivative{0}{r} \right) +
        \partialderivative{}{y} \left( C_{11}(y)\partialderivative{N_{12}(y)}{y}\frac{u^{(0)}}{r} \right) = -\partialderivative{C_{11}}{y}\DispDerivative{0}{r} - \partialderivative{C_{12}}{y}\frac{u_r^{(0)}}{r},\\
       \DispDerivative{0}{r} \partialderivative{}{y} \left( C_{11}(y)\partialderivative{N_{11}(y)}{y} \right) +
       \frac{u_r^{(0)}}{r} \partialderivative{}{y} \left( C_{11}(y)\partialderivative{N_{12}(y)}{y}(y) \right) = 
       -\partialderivative{C_{11}}{y}\DispDerivative{0}{r} - \partialderivative{C_{12}}{y}\frac{u_r^{(0)}}{r}.
    \end{aligned}
\end{equation}

Comparing terms in Eq. (\ref{eq:PreLocalProblems}) by its displacement derivative coefficient, the following holds:

\begin{equation}
    \begin{aligned}
        \partialderivative{}{y} \left( C_{11}\partialderivative{N_{11}(y)}{y} \right) = -\partialderivative{C_{11}}{y}, \\
        \partialderivative{}{y} \left( C_{11}\partialderivative{N_{12}(y)}{y} \right) = -\partialderivative{C_{12}}{y}.
    \end{aligned}
\end{equation}

Rearranging,

\begin{equation} \label{eq:LocalProblems}
    \begin{aligned}
        \partialderivative{}{y} \left( C_{11}\partialderivative{N_{11}(y)}{y} + C_{11}\right) = 0, \\
        \partialderivative{}{y} \left( C_{11}\partialderivative{N_{12}(y)}{y} + C_{12}\right) = 0.
    \end{aligned}
\end{equation}

\RaggedRight The result found in Eq. (\ref{eq:LocalProblems}) is the \textbf{local problems}, governing the unit cell's behavior. To ensure the existence and uniqueness of the local problem, the following boundary conditions must apply, according to \ref{}.

\begin{equation} \label{eq:LocalProblemsWithBCs}
    \begin{aligned}
        \partialderivative{}{y} \left( C_{11}\partialderivative{N_{11}(y)}{y} + C_{11}\right) = 0, \\
        \partialderivative{}{y} \left( C_{11}\partialderivative{N_{12}(y)}{y} + C_{12}\right) = 0. \\
         %\text{Interface conditions:} \\
        [[N_{11}(y)]] = [[N_{12}(y)]] = 0, \quad y\in {y_I} && \quad(\text{Interface condition}).\\
         %\text{Uniqueness condition:}\\
        \langle N^{(1)}_{11}(y) \rangle = \langle N^{(1)}_{12}(y) \rangle = 0, \quad y \in [0, y_I],  \\ 
        \langle N^{(2)}_{12}(y) \rangle = \langle N^{(2)}_{12}(y) \rangle =0, \quad y \in [ y_I, y_e ],  && (\text{Uniqueness conditions}).
    \end{aligned}
\end{equation}

where $y_I$ is the interface position and $y_e$ the external surface of the unit cell. Integrating both expressions with respect to y yields:

\begin{equation} \label{eq:IntegratedLocalProblems}
    \begin{aligned}
         C_{11}(y)\partialderivative{N_{11}(y)}{y} + C_{11}(y) = K_1, \\
         C_{11}(y)\left( 1 + \partialderivative{N_{11}(y)}{y} \right) = K_1, \\
        \partialderivative{N_{11}(y)}{y} = \frac{K_1}{C_{11}(y)} - 1.\\
         C_{11}(y)\partialderivative{N_{12}(y)}{y} + C_{12}(y) = K_2, \\
         \partialderivative{N_{12}(y)}{y} = \frac{K_2 - C_{12}(y)}{C_{11}(y)}.
    \end{aligned}
\end{equation}

Applying the average operator from \ref{eq:LocalProblemsWithBCs}:

\begin{equation}
    \begin{aligned}
        \left\langle\partialderivative{N_{11}(y)}{y}\right\rangle = \frac{1}{Y}\int_Y\partialderivative{N_{11}(y)}{y} dY = ... = 0 = \\
        \left\langle \frac{K_1}{C_{11}(y)} - 1 \right\rangle.
    \end{aligned}
\end{equation}

The expression for $u_r^{(1)}$ can also be used in Eq. (\ref{eq:HookesLawWithSAFRdirection}), yielding:

\begin{equation}
    \begin{aligned} \label{eq:LocalProblemsinSAFStress}
       &\sigma^{(0)}_r =  C_{11}(y) \partialderivative{u_r^{(0)}}{r} + C_{11}(y) \partialderivative{u_r^{(1)}}{y} + C_{12}(y)\DeformationTheta,\\
       &\sigma^{(0)}_r =  C_{11}(y) \partialderivative{u_r^{(0)}}{r} + C_{12}(y)\DeformationTheta + C_{11}(y) \partialderivative{}{y}\left( 
        \LocalFunctions
       \right),\\
        &\sigma^{(0)}_r = C_{11}(y) \partialderivative{u_r^{(0)}}{r} + C_{12}(y)\DeformationTheta + C_{11}(y) \partialderivative{N_{11}(y)}{y}\DispDerivative{0}{r} + C_{11}(y)\partialderivative{N_{12}(y)}{y}\DeformationTheta,\\
        &\sigma^{(0)}_r = \left( C_{11}(y) + C_{11}(y) \partialderivative{N_{11}(y)}{y}\right)\partialderivative{u_r^{(0)}}{r} + \left(C_{12}(y)\ + C_{11}(y)\partialderivative{N_{12}(y)}{y}\right)\DeformationTheta.\\
    \end{aligned}
\end{equation}

Applying the average in Eq. (\ref{eq:LocalProblemsinSAFStress}), it becomes:

\begin{equation}
    \begin{aligned}\label{eq:EffectiveHookesLawR}
        &\hat{\sigma}^{(0)}_r(r) = \langle \sigma^{(0)}_r(r,y)\rangle = \\
        &\left \langle C_{11}(y) + C_{11}(y) \partialderivative{N_{11}(y)}{y}\right\rangle \partialderivative{u_r^{(0)}}{r} + \left \langle C_{12}(y)\ + C_{11}(y)\partialderivative{N_{12}(y)}{y}\right \rangle\DeformationTheta \\
          %\hat{C}_{11} \partialderivative{u_r^{(0)}}{r} + \hat{C}_{12}\DeformationTheta.
    \end{aligned}
\end{equation}

Comparing Eq. (\ref{eq:EffectiveHookesLawR}) to Hooke's law in (\ref{eq:HookesLawCylindricalCoords}), it can be observed that the zero order expansion for $\sigma_r^{\varepsilon}(r)$ has constant versions for the elastic coefficients. These coefficients are called \textcolor{blue}{effective coefficients}. The same can be found for $\sigma_\theta^{\varepsilon}(r)$ by returning to Eq. (\ref{eq:HookesLawWithSAF_MatrixForm}). By the symmetry of the simplified elastic matrix, $\sigma_\theta^{\varepsilon}(r)$ will be a version of $\sigma_r^{\varepsilon}(r)$ as written in Eq. (\ref{eq:SAFAndHookeInRDirection}), only with switched coefficients as such:

\begin{equation}
    \begin{aligned} \label{eq:HookesLawWithSAFThetaDirection}
        &\varepsilon^{-1}:\partialderivative{u_r^{(0)}}{y} = 0,\\
        &\varepsilon^0:\sigma^{(0)}_\theta = C_{12}(y) \partialderivative{u_r^{(0)}}{r} + C_{12}(y) \partialderivative{u_r^{(1)}}{y} + C_{11}(y)\DeformationTheta, \\
        &\varepsilon^1: \sigma^{(1)}_\theta = C_{12}(y)\partialderivative{u_r^{(1)}}{r} + C_{12}(y)\partialderivative{u_r^{(2)}}{y}+ C_{11}(y)\frac{u_r^{(1)}}{r} , \\
        &\varepsilon^2: \sigma^{(2)}_\theta = C_{12}(y)\partialderivative{u_r^{(2)}}{r} +C_{11}(y) \frac{u_r^{(2)}}{r}.
    \end{aligned}
\end{equation}

Considering that $u^{(1)}(r)$ maintains the form $ u_r^{(1)} = N_{11}(y) \DispDerivative{0}{r} + N_{12}(y)\frac{u_r^{(0)}}{r} $, the zero-order term in (\ref{eq:HookesLawWithSAFThetaDirection}) assumes the following form:

\begin{equation}
    \begin{aligned} \label{eq:LocalProblemsInSAFStressTheta}
        &\sigma^{(0)}_\theta =  C_{12}(y) \partialderivative{u_r^{(0)}}{r} + C_{12}(y) \partialderivative{u_r^{(1)}}{y} + C_{11}(y)\DeformationTheta,\\
       &\sigma^{(0)}_\theta =  C_{12}(y) \partialderivative{u_r^{(0)}}{r} + C_{11}(y)\DeformationTheta + C_{12}(y) \partialderivative{}{y}\left( 
        \LocalFunctions
       \right),\\
        &\sigma^{(0)}_\theta = C_{12}(y) \partialderivative{u_r^{(0)}}{r} + C_{11}(y)\DeformationTheta + C_{12}(y) \partialderivative{N_{11}(y)}{y}\DispDerivative{0}{r} + C_{12}(y)\partialderivative{N_{12}(y)}{y}\DeformationTheta,\\
        &\sigma^{(0)}_\theta = \left( C_{12}(y) + C_{12}(y) \partialderivative{N_{11}(y)}{y}\right)\partialderivative{u_r^{(0)}}{r} + \left(C_{11}(y)\ + C_{12}(y)\partialderivative{N_{12}(y)}{y}\right)\DeformationTheta.\\
    \end{aligned}
\end{equation}

Applying the average in Eq. (\ref{eq:LocalProblemsInSAFStressTheta}), it becomes:

\begin{equation}
    \begin{aligned}\label{eq:EffectiveHookesLawTheta}
        &\hat{\sigma}^{(0)}_\theta(r) = \langle \sigma^{(0)}_\theta(r,y)\rangle = \\
        &\left \langle C_{12}(y) + C_{12}(y) \partialderivative{N_{11}(y)}{y}\right\rangle \partialderivative{u_r^{(0)}}{r} + \left \langle C_{11}(y)\ + C_{12}(y)\partialderivative{N_{12}(y)}{y}\right \rangle\DeformationTheta. \\
          %\hat{C}_{11} \partialderivative{u_r^{(0)}}{r} + \hat{C}_{12}\DeformationTheta.
    \end{aligned}
\end{equation}

Eqs. (\ref{eq:EffectiveHookesLawR}) and (\ref{eq:EffectiveHookesLawTheta}) are the effective stresses. Now, the zero-order equation in (\ref{eq:EquilibriumWithStressSAFs}) can be developed:

\begin{equation}
    \begin{aligned} \label{eq:EquilibriumWithStressSAFs}
            \partialderivative{\sigma_r^{(0)}(r)}{r} + \partialderivative{\sigma_r^{(1)}(r,y)}{y} + \frac{1}{r} \left(  \sigma_r^{(0)}(r) - \sigma_{\theta}^{(0)}(r,y) \right) = 0.
        \end{aligned}
\end{equation}

Averaging Eq. (\ref{eq:EquilibriumWithStressSAFs}) gives:

\begin{equation}
    \begin{aligned} \label{eq:EquilibriumWithStressSAFsAveraged}
            \left\langle \partialderivative{\sigma_r^{(0)}(r)}{r} + \partialderivative{\sigma_r^{(1)}(r,y)}{y} + \frac{1}{r} \left(  \sigma_r^{(0)}(r) - \sigma_{\theta}^{(0)}(r,y) \right) \right\rangle = \\
            \left\langle \partialderivative{\sigma_r^{(0)}(r)}{r}\right\rangle + \left\langle\partialderivative{\sigma_r^{(1)}(r,y)}{y}\right\rangle + \left\langle\frac{1}{r} \left( \sigma_r^{(0)}(r) - \sigma_{\theta}^{(0)}(r,y) \right) \right\rangle.
        \end{aligned}
\end{equation}

Considering the assumption that $\sigma^{(1)}(r,y)$ is $y$-periodic, the average of its derivative will be 0. This can be verified by applying the product rule to the average of the derivative of $\sigma^{(1)}(r,y)$, as such:

\begin{equation}
    \begin{aligned}
        &\left\langle\partialderivative{\sigma_r^{(1)}(r,y)}{y}\right\rangle =\\ &\frac{1}{Y}\int_0^{y_e} \partialderivative{\sigma_r^{(1)}(r,y)}{y} J(y)dy =\\&
        \frac{1}{Y}\left[ (\sigma_r^{(1)}(r,y)|_{y=y_e} - \sigma_r^{(1)}(r,y)|_{y=0}) J(y) - \int_0^{y_e}(\sigma_r^{(1)}(r,y)|_{y=y_e} - \sigma_r^{(1)}(r,y)|_{y=0})\partialderivative{J(y)}{y}dy\right] = 0,
    \end{aligned}
\end{equation}

where $J(y)$ is the \textcolor{blue}{jacobian} of the curvilinear system and $y_e$ the outer radius of the unit cell. Then, Eq. (\ref{eq:EquilibriumWithStressSAFsAveraged}) becomes:

\begin{equation}
    \begin{aligned}\label{eq:AlmostHomogenizedEquation}
        &\left\langle \partialderivative{\sigma_r^{(0)}(r)}{r}\right\rangle + \frac{1}{r}\left\langle  \sigma_r^{(0)}(r) - \sigma_{\theta}^{(0)}(r,y)  \right\rangle = \\
        &\left\langle \partialderivative{}{r} \left( \StressROrderZero \right)\right\rangle + \frac{1}{r}\left\langle \StressROrderZero - \left( \StressThetaOrderZero \right)  \right\rangle.
    \end{aligned}
\end{equation}

Now all functions in (\ref{eq:AlmostHomogenizedEquation}) are functions of $r$ only, so it can be simplified to:

\begin{equation}
    \begin{aligned}\label{eq:HomogenizedEquation}
        & \partialderivative{}{r} \left( \StressROrderZero \right) + 
        \frac{1}{r} \hat{C}_{11} \partialderivative{u_r^{(0)}}{r} + 
        \hat{C}_{12}\frac{u_r^{(0)}}{r^2} - 
        \frac{1}{r}\left( \StressThetaOrderZero \right) = 0, \\
        & \hat{C}_{11}\partialderivative{^2 u_r^{(0)}}{r^2} + 
        \hat{C}_{12}\partialderivative{}{r}\left( \frac{u_r^{(0)}}{r} \right) + 
        \frac{1}{r} \partialderivative{u_r^{(0)}}{r}(\hat{C}_{11} - \hat{C}_{12}) +
        \frac{u_r^{(0)}}{r^2}(\hat{C}_{12} - \hat{C}_{11}),\\
        & \hat{C}_{11}\partialderivative{^2 u_r^{(0)}}{r^2} + 
        \hat{C}_{12}\left( \frac1r \partialderivative{u_r^{(0)}}{r} - \frac{u_r^{(0)}}{r^2} \right) + 
        \frac{1}{r} \partialderivative{u_r^{(0)}}{r}(\hat{C}_{11} - \hat{C}_{12}) +
        \frac{u_r^{(0)}}{r^2}(\hat{C}_{12} - \hat{C}_{11}),\\
        & \hat{C}_{11}\partialderivative{^2 u_r^{(0)}}{r^2} + 
        \hat{C}_{11}\frac{1}{r} \partialderivative{u_r^{(0)}}{r} -
        \hat{C}_{11}\frac{u_r^{(0)}}{r^2}.
    \end{aligned}
\end{equation}

Which has the same structure as Eq. (\ref{eq:EquilibriumCylindricalForm}), but now with effective coefficients replacing the fast oscillating ones. For that reason, Eq. (\ref{eq:HomogenizedEquation}), is called the \textcolor{blue}{Homogenized Equation}.